% =====================================================================
%  bg_B_language_models.tex
%  WHERE:  inside Section 2.3 "Language Models",
%          directly AFTER your subsection 2.3.3 "Fine-Tuning".
%  ADDS:   2.3.4 Parameter-Efficient Fine-Tuning (LoRA / QLoRA)
%          2.3.5 Prompting, In-Context Learning and Chain-of-Thought
%          2.3.6 Hallucination in Large Language Models
% =====================================================================

\subsection{Parameter-Efficient Fine-Tuning}
\label{sec:bg-peft}

Full fine-tuning updates \emph{all} parameters of a model. For current LLMs
with billions of parameters, this requires a large amount of GPU memory and
produces a complete copy of the model for every task. \emph{Parameter-efficient
fine-tuning} (PEFT) methods avoid this by freezing the pre-trained weights and
training only a small number of additional parameters. The most widely used
PEFT method is \emph{Low-Rank Adaptation} (LoRA)~\cite{hu2022lora}.

LoRA is based on the observation that the weight change needed to adapt a
pre-trained model to a new task has a low ``intrinsic rank''. Instead of
learning a full update matrix, LoRA learns two small matrices whose product
approximates it (Figure~\ref{fig:bg-lora}).

\begin{bgdefinition}{Low-Rank Adaptation (LoRA)}
Let $W_0 \in \mathbb{R}^{d \times k}$ be a frozen pre-trained weight matrix
of a linear layer. LoRA represents the task-specific update as
$\Delta W = BA$ with $B \in \mathbb{R}^{d \times r}$,
$A \in \mathbb{R}^{r \times k}$ and rank $r \ll \min(d,k)$. For an input
$x$, the adapted layer computes
\[
  h \;=\; W_0 x + \frac{\alpha}{r}\, B A x ,
\]
where $\alpha$ is a constant scaling factor. Only $A$ and $B$ are trained;
$B$ is initialised with zeros so that training starts from the unchanged
pre-trained model~\cite{hu2022lora}.
\end{bgdefinition}

\begin{figure}[htbp]
  \centering
  \begin{tikzpicture}[
      font=\small,
      >={Stealth[length=2.2mm]},
      mat/.style={draw, thick, minimum height=2.6cm, align=center},
      lab/.style={font=\footnotesize, align=center}]
    % input
    \node[draw, thick, minimum width=2.6cm, minimum height=0.55cm, fill=bgGreyLight] (x) at (0,0) {input $x \in \mathbb{R}^{k}$};
    % frozen W0
    \node[mat, minimum width=2.6cm, fill=bgBlueLight, draw=bgBlue] (W) at (-2.1,2.2)
        {\textbf{Pre-trained}\\\textbf{weights} $W_0$\\[2pt]$d \times k$\\[2pt]\footnotesize(frozen)};
    % LoRA A and B
    \node[draw=bgOrange, thick, fill=bgOrangeLight, trapezium, trapezium angle=65,
          minimum width=2.3cm, minimum height=0.9cm, shape border rotate=180] (A) at (2.1,1.45) {$A$ \;\footnotesize($r\times k$)};
    \node[draw=bgOrange, thick, fill=bgOrangeLight, trapezium, trapezium angle=65,
          minimum width=2.3cm, minimum height=0.9cm] (B) at (2.1,2.95) {$B$ \;\footnotesize($d\times r$)};
    % sum
    \node[draw, thick, circle, inner sep=1pt] (plus) at (0,4.55) {$+$};
    \node[draw, thick, minimum width=2.6cm, minimum height=0.55cm, fill=bgGreyLight] (h) at (0,5.5) {output $h \in \mathbb{R}^{d}$};
    % arrows
    \draw[->, thick] (x.north) -- ++(0,0.35) -| (W.south);
    \draw[->, thick] (x.north) -- ++(0,0.35) -| (A.south);
    \draw[->, thick] (A.north) -- (B.south) node[midway, right, lab] {\ rank $r$};
    \draw[->, thick] (W.north) |- (plus.west);
    \draw[->, thick] (B.north) |- (plus.east) node[pos=0.25, right, lab] {$\times\frac{\alpha}{r}$};
    \draw[->, thick] (plus) -- (h);
    % annotations
    \node[lab, text=bgBlue, anchor=east] at ($(W.west)+(-0.1,0)$) {not\\updated};
    \node[lab, text=bgOrange, anchor=west] at ($(B.east)!0.5!(A.east)+(0.15,0)$) {trainable\\LoRA adapter};
  \end{tikzpicture}
  \caption[Low-Rank Adaptation (LoRA)]{Low-Rank Adaptation (LoRA). The
  pre-trained weight matrix $W_0$ stays frozen; only the two small matrices
  $A$ and $B$ are trained, and their scaled product is added to the output of
  the original layer. Adapted from Hu et al.~\cite{hu2022lora}.}
  \label{fig:bg-lora}
\end{figure}

\begin{bgexample}[How many parameters does LoRA train?]
Consider a single attention projection of an 8-billion-parameter model with
$d = k = 4096$. Full fine-tuning updates
$d \cdot k = 4096 \cdot 4096 = 16{,}777{,}216$ parameters in this matrix.
With LoRA of rank $r=16$, only
$r(d+k) = 16 \cdot (4096 + 4096) = 131{,}072$ parameters are trained, which
is about $0.78\,\%$ of the original matrix. The resulting adapter can be
stored as a small separate file and loaded on top of the unchanged base
model.
\end{bgexample}

\emph{QLoRA}~\cite{dettmers2023qlora} further reduces memory requirements by
storing the frozen base model in 4-bit precision while training the LoRA
matrices in higher precision. This combination makes it possible to fine-tune
models with billions of parameters, such as the Llama~3
family~\cite{grattafiori2024llama3}, on a single GPU. Because the base model
remains unchanged, several adapters, for example one trained with and one
trained without knowledge-graph rules, can be compared on top of exactly the
same model, which is the experimental setting used later in this thesis.

\subsection{Prompting, In-Context Learning and Chain-of-Thought}
\label{sec:bg-prompting}

Besides changing the parameters of a model, its behaviour can be steered
through its \emph{input}. A \emph{prompt} is the text that is given to an LLM
to elicit a response; it typically combines an instruction, optional
context (e.g.\ facts retrieved from a Knowledge Graph) and the actual
question. Brown et al.~\cite{brown2020language} showed that sufficiently
large models can solve new tasks from a few demonstrations placed in the
prompt, without any parameter update. This ability is called
\emph{in-context learning}, and a prompt that contains demonstrations is
called a \emph{few-shot} prompt.

For tasks that require several reasoning steps, it helps to ask the model to
write down its intermediate steps before answering.

\begin{bgdefinition}{Chain-of-Thought (CoT)}
A \emph{chain of thought} is a sequence of intermediate natural-language
reasoning steps $s_1, \dots, s_m$ that leads from an input question $q$ to a
final answer $a$. \emph{Chain-of-thought prompting} provides demonstrations
of the form $(q, s_1, \dots, s_m, a)$ so that the model generates its own
reasoning steps before producing the answer~\cite{wei2022chain}.
\end{bgdefinition}

Chain-of-thought demonstrations can be used in two ways: at inference time
inside the prompt, as originally proposed by Wei et al.~\cite{wei2022chain},
or as \emph{training data} for supervised fine-tuning, so that the model
learns to produce explicit reasoning by itself. The second option is the one
followed in this thesis, where the reasoning steps are derived from rules
over a Knowledge Graph (Section~\ref{sec:bg-rules-to-cot}).

\begin{bgexample}[Standard answer vs.\ chain-of-thought answer]
\textbf{Question:} \emph{``My ejection fraction is 25\,\%. Which type of
heart failure do I have?''}\par\smallskip
\textbf{Standard answer:} \emph{``You have heart failure with reduced ejection
fraction.''}\par\smallskip
\textbf{Chain-of-thought answer:}
\begin{enumerate}
  \item The left-ventricular ejection fraction (LVEF) describes how much blood
        the left ventricle pumps out with each beat.
  \item The ESC guideline classifies heart failure with an LVEF of
        $\leq 40\,\%$ as heart failure with reduced ejection fraction
        (HFrEF)~\cite{mcdonagh2021esc}.
  \item $25\,\% \leq 40\,\%$.
  \item \emph{Therefore}, an LVEF of 25\,\% corresponds to HFrEF. The exact
        diagnosis is made by the treating cardiologist.
\end{enumerate}
Both answers have the same conclusion, but only the second one makes the
underlying medical knowledge and the reasoning step visible, so that a
clinician can check each step.
\end{bgexample}

\subsection{Hallucination in Large Language Models}
\label{sec:bg-hallucination}

LLMs generate text by repeatedly predicting likely next tokens. Fluency is
therefore not a guarantee of correctness: a model can produce statements that
sound convincing but are not supported by any source. This phenomenon is
known as \emph{hallucination}~\cite{ji2023survey}.

\begin{bgdefinition}{Hallucination}
A \emph{hallucination} is generated content that is nonsensical or unfaithful
to the provided source content or to established world knowledge. It is
\emph{intrinsic} if the output contradicts the source, and \emph{extrinsic}
if the output cannot be verified from the source, i.e.\ it can neither be
supported nor contradicted by it~\cite{ji2023survey}.
\end{bgdefinition}

\begin{bgexample}[A hallucination in patient education]
A patient asks: \emph{``My heart is functioning at 25\,\%. Am I eligible for
a heart transplant?''} A model answers: \emph{``Yes. Every patient with an
ejection fraction below 30\,\% qualifies for a heart transplant.''} The
answer is fluent and contains a plausible-looking threshold, but no
guideline states such a rule; eligibility depends on many additional clinical
factors~\cite{mcdonagh2021esc}. If the prompt contained a guideline excerpt
describing the actual eligibility criteria, the answer contradicts its source
and is an \emph{intrinsic} hallucination; if no such source was given, the
claim is unsupported and cannot be verified, which makes it an
\emph{extrinsic} hallucination.
\end{bgexample}

In medicine, hallucinations are especially critical, because patients may
act on incorrect information and often lack the background to detect
errors~\cite{thirunavukarasu2023large}. Two complementary strategies are
used to reduce them: (i)~\emph{grounding} the model in curated, explicit
knowledge, such as a Knowledge Graph built from clinical guidelines
(Section~\ref{sec:bg-kg-llm}), and (ii)~making the model's reasoning
explicit, so that each step can be checked against that knowledge
(Section~\ref{sec:bg-nesy}). This thesis combines both strategies.
